\documentclass[10pt]{article}
\usepackage[T1]{fontenc}
\usepackage[utf8]{inputenc}
\usepackage[french]{babel}
\usepackage{amsmath,amssymb,mathrsfs,mathpazo}
\usepackage{geometry}
\geometry{a4paper,margin=2cm}
\usepackage{graphicx,array,tabularx,tabularray,colortbl,multirow,multicol,booktabs,enumitem,pifont,ulem,fancybox,diagbox,cancel}
\usepackage{pstricks,pstricks-add,pst-plot,pst-tree,pst-node,pst-text,pst-func,pst-math,pst-bspline,pst-3dplot,pst-eucl,pst-solides3d}
\let\pscalculate\undefined
\usepackage{pst-bezier}
\usepackage{tikz,pgfplots,tkz-tab}
\usetikzlibrary{arrows,arrows.meta,patterns,calc,positioning,trees,decorations.markings,shapes.geometric,shapes.misc,shapes.arrows,intersections,decorations.pathmorphing,decorations.pathreplacing,backgrounds,shadows,babel}
\pgfplotsset{compat=1.18}
\usepgfplotslibrary{fillbetween}
\usepackage[np]{numprint}
\usepackage{listings,makecell,pas-tableur,twoopt}
\let\barre\undefined
\usepackage[pstricks]{tablvar}
\definecolor{darkgreen}{RGB}{0,100,0}
\providecommand{\red}{\color{red}}
\providecommand{\blue}{\color{blue}}
\providecommand{\green}{\color{green}}
\providecommand{\black}{\color{black}}
\providecommand{\white}{\color{white}}
\providecommand{\gray}{\color{gray}}
\providecommand{\cyan}{\color{cyan}}
\providecommand{\magenta}{\color{magenta}}
\providecommand{\yellow}{\color{yellow}}
\providecommand{\violet}{\color{violet}}
\providecommand{\orange}{\color{orange}}
\providecommand{\euro}{\texteuro}
\providecommand{\EUR}{\texteuro}
\providecommand{\ieme}{\textsuperscript{e}}
\providecommand{\degres}{\ensuremath{^\circ}}
\providecommand{\degree}{\ensuremath{^\circ}}
\providecommand{\vv}[1]{\overrightarrow{#1}}
\providecommand{\Coord}[1]{#1}
\providecommand{\decosix}{}
\providecommand{\starredbullet}{\ensuremath{\bullet}}
\pagestyle{empty}
\setlength{\parindent}{0pt}
\usepackage{adjustbox}\DeclareUnicodeCharacter{2236}{:}\DeclareUnicodeCharacter{2190}{\ensuremath{\leftarrow}}
\providecommand{\R}{\mathbb{R}}
\providecommand{\N}{\mathbb{N}}
\providecommand{\Z}{\mathbb{Z}}
\providecommand{\Q}{\mathbb{Q}}
\providecommand{\D}{\mathbb{D}}
\providecommand{\C}{\mathbb{C}}
\providecommand{\e}{{\mathrm{e}}}
\providecommand{\vect}[1]{\overrightarrow{#1}}
\providecommand{\vectt}[1]{\overrightarrow{\mathrm{#1}}}
\providecommand{\vvt}[1]{\overrightarrow{\mathrm{#1}}}
\providecommand{\np}[2][]{{#2~#1}}
\providecommand{\vv}[1]{{\overrightarrow{#1}}}
\renewcommand{\d}{{\mathrm{d}}}
\providecommand{\degres}{{^\circ}}
\providecommand{\degre}{{^\circ}}
\providecommand{\blue}{\color{blue}}
\providecommand{\red}{\color{red}}
\providecommand{\enskip}{\quad}
\providecommand{\llbracket}{\lbrack\!\lbrack}
\providecommand{\rrbracket}{\rbrack\!\rbrack}
\providecommand{\euro}{{\text{€}}}
\providecommand{\no}{{\text{n°}}}
\providecommand{\interval}[2]{{[#1;#2]}}
\providecommand{\intervFO}[2]{{[#1;#2[}}
\providecommand{\dotfill}{\ldots}

\providecommand{\ds}{\displaystyle}
\providecommand{\textding}[1]{#1}
\providecommand{\Oij}{\(\mathrm{O};\vec{\imath},\vec{\jmath}\)}
\providecommand{\Oijk}{\(\mathrm{O};\vec{\imath},\vec{\jmath},\vec{k}\)}
\providecommand{\Ouv}{\(\mathrm{O};\vec{u},\vec{v}\)}
\providecommand{\intervFF}[2]{\ensuremath{[#1;#2]}}
\renewcommand{\euro}{\ifmmode\text{\texteuro}\else\texteuro\fi}
\DeclareUnicodeCharacter{00B0}{\ifmmode^{\circ}\else\textdegree\fi}
\DeclareUnicodeCharacter{20AC}{\ifmmode\text{\texteuro}\else\texteuro\fi}
% Énoncés : Éducation nationale. Transcription et corrigés : APMEP, Denis Vergès et contributeurs. Découpage et classement : CoopMaths, Sébastien Lozano, Éric Elter et contributeurs. Adaptation de la présentation : Maths Explorer.
% CC BY-SA : https://coopmaths.fr/www/about/
% https://coopmaths.fr/alea/static/bac/2026/tex/bac_2026_06_sujet2_asie_1.tex
\begin{document}
\noindent\textit{APMEP / CoopMaths. CC BY-SA. Présentation adaptée par Maths Explorer.}\par\medskip
\subsection*{Partie  A}
On considère la fonction $f$ définie sur $\left]-\infty ; \dfrac{3}{2}\right [$ par : \[f(x)=\dfrac{x-2}{2 x-3}\]

	\begin{enumerate}
		\item Justifier tous les éléments du tableau de variation ci-dessous.
	\end{enumerate}

	\begin{center}
	\begin{tikzpicture}
		\tkzTabInit[lgt=1,espcl=5]{$x$/1, $f$/2}{$-\infty$,$\dfrac{3}{2}$}
		\tkzTabVar{-/$\dfrac{1}{2}$,+D/$+\infty$}
	\end{tikzpicture}
	\end{center}

	\begin{enumerate}
		\setcounter{enumi}{1}
		\item En déduire que pour tout $x \in[0 ; 1]$, on a :
		\[f(x) \in[0 ; 1].\]
	\end{enumerate}

\subsection*{Partie  B}
On considère la suite $(u_{n})$ définie par :
\[\begin{cases}
	u_{0}=0 \\
	u_{n+1}=\dfrac{u_{n}-2}{2 u_{n}-3}\, ,\quad \text { pour tout entier naturel } n
\end{cases}\]

	\begin{enumerate}
		\item En utilisant la fonction $f$, démontrer par récurrence que pour tout entier naturel $n$ :
		\[0 \leqslant  u_{n} \leqslant  u_{n+1} \leqslant  1.\]

		\item En déduire que la suite $(u_{n})$ converge.

		\item On note $\ell$ la limite de la suite $(u_{n})$.

		En admettant que $\ell$ est solution de l'équation $f(x)=x$ sur l'intervalle $[0 ; 1]$, montrer que $\ell=1$.

		\item On donne ci-dessous une fonction \texttt{seuil} écrite en langage Python.

		\begin{center}
		\begin{tblr}{colspec={X[6cm,font=\ttfamily]},vlines,hline{1,Z}={solid},rowsep=0pt}
			def seuil(h):\\
			\quad n = 0\\
			\quad u = 0\\
			\quad while u < 1-h:\\
			\quad \quad n = n+1\\
			\quad \quad u = (u-2)/(2*u-3)\\
			\quad return n\\
		\end{tblr}
		\end{center}

	L'appel ~\texttt{seuil(0.0001)}~ renvoie la valeur \np{5000}.

	Interpréter cette valeur dans le contexte de l'exercice.

	\item \begin{enumerate}
		\item Donner les quatre premiers termes de la suite $\left(u_{n}\right)$ sous forme de fractions irréductibles.

		\item Conjecturer l'expression de $u_{n}$ en fonction de $n$ et démontrer cette conjecture.
	\end{enumerate}
	\end{enumerate}

\end{document}
